\subsection{Тестирование}

\subsubsection{Автоматизированное тестирование серверной части}

Для обеспечения качества серверной части разработаны автоматизированные тесты на платформе xUnit. Модульное тестирование охватывает классы PasswordServiceTests и AuthServiceSignInTests (хеширование пароля, вход жителя, отказ для неподходящей роли); интеграционные тесты AuthEndpointsTests проверяют эндпоинты входа и регистрации, метод получения текущего пользователя и создание обращения с bearer-токеном.

На рисунке~\ref{fig:tests-backend} показан результат успешного выполнения тестов серверной части.

\begin{figure}[H]
    \centering
    \setlength{\fboxrule}{0.4pt}
    \setlength{\fboxsep}{0pt}
    \fbox{\includegraphics[width=0.85\linewidth,height=0.45\textheight,keepaspectratio]{tests-runner-result.png}}
    \caption{Результат выполнения автоматизированных тестов серверной части}
    \label{fig:tests-backend}
\end{figure}

Все тесты завершились успешно.

\subsubsection{Модульное тестирование мобильного приложения}

Для мобильного приложения разработаны модульные тесты на базе фреймворка Vitest. Общее количество тестовых файлов составляет 20 единиц. Тестирование охватывает следующие модули:

\textbf{Модуль аутентификации} (6 тестовых файлов):
\begin{itemize}
    \item auth-api.test.ts~--- проверка корректности HTTP-запросов к эндпоинтам аутентификации (/Auth/sign-in, /Auth/sign-out, /Auth/google/sign-in, /Auth/me, /Auth/forgot-password, /Auth/reset-password) и маппинг серверного ответа;
    \item auth-validation.test.ts~--- проверка валидации полей форм входа, регистрации, восстановления и сброса пароля;
    \item auth-session-manager.test.ts~--- проверка управления жизненным циклом сессии: сохранение и восстановление токенов из локального хранилища;
    \item google-sign-in-stub.test.ts, lazy-google-auth-client.test.ts~--- проверка интеграции с механизмом входа через Google;
    \item google-brand-icon.test.ts~--- проверка корректности отображения фирменной иконки Google.
\end{itemize}

\textbf{Модуль интерактивной карты} (6 тестовых файлов):
\begin{itemize}
    \item map-api.test.ts~--- проверка загрузки обращений и событий с сервера, маппинг DTO-объектов;
    \item map-model.test.ts~--- проверка логики построения ленты карты, фильтрации по слоям и категориям, группировки маркеров в стековые маркеры, формирования легенды;
    \item map-helpers.test.ts~--- проверка вспомогательных функций для работы с координатами;
    \item map-search.test.ts, map-search-history-state.test.ts~--- проверка функционала поиска и сохранения истории поисковых запросов;
    \item request-guard.test.ts~--- проверка механизма предотвращения дублирования одновременных сетевых запросов.
\end{itemize}

\textbf{Модуль обращений} (2 тестовых файла):
\begin{itemize}
    \item reports-api.test.ts~--- проверка операций с обращениями: получение списка с фильтрацией по статусу, получение по идентификатору, создание нового обращения, поиск дубликатов, подтверждение существующего обращения;
    \item report-errors.test.ts~--- проверка классификации ошибок API по HTTP-кодам (401, 403, 404, 400, 422, 500) и формирование пользовательских сообщений.
\end{itemize}

\textbf{Модуль уведомлений} (2 тестовых файла):
\begin{itemize}
    \item notifications-api.test.ts~--- проверка загрузки уведомлений с сервера;
    \item notifications-formatting.test.ts~--- проверка форматирования текста уведомлений.
\end{itemize}

\textbf{Ядро приложения} (4 тестовых файла):
\begin{itemize}
    \item install-auth-session-bridge.test.ts~--- проверка перехватчиков HTTP-клиента Axios: автоматическое добавление заголовка авторизации, пропуск авторизации для публичных эндпоинтов, автоматическое обновление токена при ответе 401;
    \item resolve-api-base-url.test.ts~--- проверка определения базового адреса API;
    \item runtime-log-buffer.test.ts, runtime-log-session.test.ts~--- проверка буферизации и сессионного управления журналом событий.
\end{itemize}

На рисунке~\ref{fig:tests-mobile} показан результат выполнения модульных тестов мобильного приложения.

\begin{figure}[H]
    \centering
    \setlength{\fboxrule}{0.4pt}
    \setlength{\fboxsep}{0pt}
    \fbox{\includegraphics[width=0.85\linewidth,height=0.45\textheight,keepaspectratio]{vittest.png}}
    \caption{Результат выполнения модульных тестов мобильного приложения}
    \label{fig:tests-mobile}
\end{figure}

Все тесты завершились успешно.

\subsubsection{Тестирование API через Swagger UI}

Для проверки корректности работы серверного программного интерфейса выполнено тестирование через встроенный интерфейс Swagger UI.

На рисунке~\ref{fig:swagger-get-reports} представлен запрос на получение списка обращений. В ответе возвращается массив записей с идентификатором, категорией, описанием, координатами и текущим статусом. Сервер возвращает код ответа 200, что подтверждает корректность чтения списка обращений.

\begin{figure}[H]
    \centering
    \setlength{\fboxrule}{0.4pt}
    \setlength{\fboxsep}{0pt}
    \fbox{\includegraphics[width=0.96\linewidth,height=0.42\textheight,keepaspectratio]{swagger-get-reports.png}}
    \caption{Запрос на получение списка обращений}
    \label{fig:swagger-get-reports}
\end{figure}

На рисунке~\ref{fig:swagger-duplicates} представлен запрос на поиск дубликатов обращений по координатам и категории. Эндпоинт GET /api/Reports/duplicates возвращает массив обращений, расположенных в радиусе 50~м от указанных координат и имеющих совпадающую категорию.

\begin{figure}[H]
    \centering
    \setlength{\fboxrule}{0.4pt}
    \setlength{\fboxsep}{0pt}
    \fbox{\includegraphics[width=0.96\linewidth,height=0.42\textheight,keepaspectratio]{swagger-find-duplicates.png}}
    \caption{Запрос на поиск дубликатов обращений}
    \label{fig:swagger-duplicates}
\end{figure}

На рисунке~\ref{fig:swagger-categories} представлен запрос на получение справочника категорий обращений. Возвращаемый массив содержит идентификатор, наименование и код каждой категории.

\begin{figure}[H]
    \centering
    \setlength{\fboxrule}{0.4pt}
    \setlength{\fboxsep}{0pt}
    \fbox{\includegraphics[width=0.96\linewidth,height=0.42\textheight,keepaspectratio]{swagger-get-categories.png}}
    \caption{Запрос на получение справочника категорий обращений}
    \label{fig:swagger-categories}
\end{figure}

На рисунке~\ref{fig:swagger-create} представлен запрос на создание нового обращения. В теле запроса передаются описание проблемы, координаты и идентификатор категории. Сервер возвращает код успешного выполнения и созданную запись.

\begin{figure}[H]
    \centering
    \setlength{\fboxrule}{0.4pt}
    \setlength{\fboxsep}{0pt}
    \fbox{\includegraphics[width=0.96\linewidth,height=0.42\textheight,keepaspectratio]{create-report.png}}
    \caption{Запрос на создание нового обращения}
    \label{fig:swagger-create}
\end{figure}

Во всех случаях система возвращала ожидаемые коды ответа и корректные структуры JSON, что подтверждает корректность работы серверного программного интерфейса.

\subsubsection{Функциональное тестирование веб-панели}

Функциональное тестирование веб-панели диспетчера включало следующие проверки:
\begin{itemize}
    \item тестирование адаптивности интерфейса на различных разрешениях экрана;
    \item проверка корректности работы интерактивной карты (масштабирование, перемещение, отображение маркеров обращений, фильтрация по статусу и категории);
    \item проверка корректности отображения дашборда, таблицы обращений и карточки обращения.
\end{itemize}

На рисунке~\ref{fig:web-adaptive} представлена проверка адаптивности интерфейса.

\begin{figure}[H]
    \centering
    \setlength{\fboxrule}{0.4pt}
    \setlength{\fboxsep}{0pt}
    \fbox{\includegraphics[width=0.96\linewidth,height=0.42\textheight,keepaspectratio]{adaptive-web-panel.png}}
    \caption{Адаптивное отображение веб-панели диспетчера}
    \label{fig:web-adaptive}
\end{figure}

На рисунке~\ref{fig:web-map-testing} представлена интерактивная карта в режиме функционального тестирования: активный фильтр по статусу, поиск по категории, маркер обращения и открытое всплывающее окно карточки. Это подтверждает корректность отображения маркеров с цветовой кодировкой, работу фильтрации и открытие подробной информации при выборе объекта на карте.

\begin{figure}[H]
    \centering
    \setlength{\fboxrule}{0.4pt}
    \setlength{\fboxsep}{0pt}
    \fbox{\includegraphics[width=0.96\linewidth,height=0.42\textheight,keepaspectratio]{testing-map-features.png}}
    \caption{Функциональное тестирование интерактивной карты обращений}
    \label{fig:web-map-testing}
\end{figure}

Карта отображается корректно: диспетчер может изменять масштаб, перемещаться по карте, применять фильтры и открывать карточки обращений непосредственно из картографического интерфейса.
